% SPDX-License-Identifier: LPPL-1.3c
% Copyright (C) 2026 Christophe Jorssen
% Author and Current Maintainer: Christophe Jorssen <christophe.jorssen@gmail.com>
% LPPL maintenance status: maintained
% This file is part of luafemm. See LICENSE and LICENSES.md for its terms.
% Original narrative and examples, following the progressive presentation of
% the PGF tutorials (a concrete task, small additions, then complete results).
\clearpage
\section{Tutorial: Two magnets for Maya}
\label{sec:tutorial}

Maya is preparing a lesson on magnetic circuits. She has a drawing of a
U-shaped iron core, a coil, and an armature above the two poles. The drawing
is quite satisfactory, except for the field lines: she has drawn them by hand,
and moving the armature means drawing them all again. She would also like to
answer a question that the drawing alone cannot settle: how strong is the
field in each air gap?

Her plan is to let the paths in her TikZ picture describe the physical model,
ask LuaTeX to compute the field, and use PGFPlots to graph measurements along
a contour. Later she will remove the coil and replace part of the iron with
alnico. Most of the drawing should survive this change.

We shall follow her through both examples. The complete, compilable sources
are \nolinkurl{examples/tutorial-coil.tex} and
\nolinkurl{examples/tutorial-alnico.tex}. They use the same catalogue material,
|pure-iron|, for the soft iron, so there is no B--H table to type. The earlier
|u-electromagnet.tex| example instead uses an illustrative iron curve; its
numerical values need not agree with this tutorial.

\subsection{A useful first decision: do less work}

Before choosing colours, Maya chooses a fairly coarse mesh. She will move
labels, adjust the drawing, and change her mind about the graph several
times. Rebuilding a very fine finite-element model after each such change
would be a poor use of her time.

Her starting value is |mesh size=5|, in millimetres. This is a target spacing,
not a requested number of triangles. Decreasing it generally increases the
number of triangles and the cost of the solve; in a fairly uniform planar
mesh, halving the spacing gives roughly four times as many elements. Geometry
constraints and local refinement change the actual count. The solution time
need not scale by the same factor.

\begin{quote}
\textbf{Draft the picture with a coarse global mesh. Refine the important
regions as needed, then refine and check the numerical result when the
geometry and presentation are nearly final.}
\end{quote}

There is one qualification: coarse must not mean blind to the physics.
Maya's air gaps are narrow. She keeps them explicitly in the geometry and
uses smaller elements there, even in a draft. A plausible-looking picture
on a coarse mesh is useful for composition; it is not yet an accurate
field measurement. More arrows, smoother plot styling, or more measurement
samples would not repair an insufficient mesh.

\subsection{Setting up the document}

Maya uses LuaLaTeX and loads two graphics libraries:
\begin{codeexample}[code only]
\documentclass{article}
\usepackage{tikz,pgfplots}
\usetikzlibrary{femm}
\usepgfplotslibrary{femm}
\pgfplotsset{compat=1.18}
\begin{document}
% The model picture, followed by a separate PGFPlots picture.
\end{document}
\end{codeexample}
Compile with |lualatex --no-shell-escape|. Nothing is sent to an external
solver. The complete example files use the |standalone| class to crop their
output; |article| is sufficient for following the steps here. Plain LuaTeX
and ConTeXt MkIV use the loading commands from Section~\ref{sec:start};
the physical keys and the plotting commands are the same.

The following code fragments are successive additions to the same document.
In particular, the physical paths go in one model picture, while the axes
will go in separate pictures after it. A graph is not a new magnetic problem.

\subsection{First, just draw the iron}

Maya begins with ordinary TikZ paths. The U is a single closed polygon, and
the armature is a rectangle above it. The top of each pole is at $y=30$ mm;
the bottom of the armature is at $y=33$ mm, leaving two 3 mm gaps.
\begin{codeexample}[width=5cm]
\begin{tikzpicture}[x=1mm,y=1mm,scale=.5]
  \filldraw[fill=gray!15]
    (-40,-30)--(40,-30)--(40,30)--(25,30)
    --(25,-15)--(-25,-15)--(-25,30)
    --(-40,30)--cycle;
  \filldraw[fill=gray!15]
    (-40,33) rectangle (40,48);
\end{tikzpicture}
\end{codeexample}

So far, this is only a drawing. The gray fill does not tell the solver that
there is iron inside it. To give the drawing that meaning, Maya introduces
a style that contains both physical and graphical options:
\begin{codeexample}[code only]
\tikzset{
  core/.style={femm/region={material=pure-iron},
    draw=black!70,fill=gray!15}
}
\end{codeexample}
She replaces |[fill=gray!15]| by |[core]| on both paths. The choice
|material=pure-iron| now selects a nonlinear material from the FEMM catalogue.
The colour remains an ordinary TikZ choice; changing it will not change the
material law.

Next she replaces the options of the enclosing picture:
\begin{codeexample}[code only]
\begin{tikzpicture}[femm/problem={xmin=-110,xmax=110,
  ymin=-95,ymax=115,mesh size=5}]
  % The two \filldraw[core] paths go here.
\end{tikzpicture}
\end{codeexample}
The |femm/problem| key installs millimetre coordinates and creates a rectangular
air domain larger than the device. Its outside boundary has $A_z=0$.
There is no need to draw that rectangle, nor to fill every unassigned part
with an air region. In fact, a large air region drawn \emph{after} the core
would replace the iron there: later physical regions take precedence.

Maya removes |scale=.5| for the working picture. A scale applied to the whole
model picture would only change its displayed size; it would not halve its
physical dimensions. By contrast, scaling a region inside the picture changes
the model geometry. Keeping this distinction in mind saves some surprisingly
unproductive experiments.

\subsection{Giving the coil a current}

In this two-dimensional section, the winding appears as two rectangles on
opposite sides of the left leg. Current leaves the page through one section
and enters it through the other. Maya first defines their common appearance
and material:
\begin{codeexample}[code only]
\tikzset{
  winding/.style={femm/region={material=copper},
    draw=orange!70!black,fill=orange!25}
}
\end{codeexample}
Then she adds two paths, inside the model picture and after the iron:
\begin{codeexample}[code only]
\filldraw[winding,femm/current density=6000/(6*26)*1000000]
  (-48,-10) rectangle (-42,16);
\filldraw[winding,femm/current density=-6000/(6*26)*1000000]
  (-23,-10) rectangle (-17,16);
\end{codeexample}
Each section is $6\times26$ mm$^2$. For $NI=6000$ ampere-turns, its imposed
current density is
\[
 J_z=\frac{NI}{S}=\frac{6000}{6\times26}\times10^6\ \mathrm{A/m^2},
\]
with the opposite sign in the return section. The factor $10^6$ converts
A/mm$^2$ to A/m$^2$. The key accepts this PGF expression directly. Copper
alone does not create this winding excitation; the signed current density does.

These are homogenised winding sections, not a turn-by-turn conductor model.
If Maya changes their areas but wants to keep $NI$ fixed, she must change
the current-density expression too. The solver does not infer that intention
from the picture.

\subsection{Looking at the mesh before asking for the field}

A 5 mm global spacing is quite coarse compared with a 3 mm gap. Maya therefore
adds two invisible air rectangles with a smaller local spacing:
\begin{codeexample}[code only]
\tikzset{
  gap mesh/.style={femm/region={material=air,mesh size=1}}
}
\end{codeexample}
\begin{codeexample}[code only]
\path[gap mesh] (-40,30) rectangle (-25,33);
\path[gap mesh] (25,30) rectangle (40,33);
\end{codeexample}
These paths both assign air and request local refinement. Here they occupy
only the gaps, so placing them after the solids is safe. An invisible path
still has a physical effect when it carries |femm/region|.

To see what she has asked the mesher to do, she can finish the picture with
\begin{codeexample}[code only]
\begin{scope}
  \clip (-58,-42) rectangle (58,57);
  \pic[femm/mesh style={black!25}] {femm mesh};
\end{scope}
\end{codeexample}
The clip keeps the distant air out of the illustration. It does not reduce
the computational domain. A mesh picture builds the triangulation but does
not solve for the field, which makes it a useful intermediate check.

All material paths must come \emph{before} this request. Once meshing has
started, it is too late to add the return conductor or a forgotten refinement
region to that model. Maya now knows where her last physical path belongs.

\subsection{The first computed field lines}

The interesting part is pleasantly short. Maya replaces the mesh picture
by a field picture, keeping the same clip:
\begin{codeexample}[code only]
\begin{scope}
  \clip (-58,-42) rectangle (58,57);
  \pic[femm/lines=17,femm/arrow spacing=26mm,
    femm/field style={teal!70!black}] {femm field};
\end{scope}
\end{codeexample}
This request meshes and solves the model if necessary, then draws contours of
$A_z$. Their arrows follow $\mathbf B$. Maya may put the mesh picture first
and the field picture second to see both; the second request reuses the mesh.
Labels and the $\odot$, $\otimes$ conductor symbols are added afterwards so
they remain readable.

The 17 in |femm/lines=17| counts potential levels, not triangles. Changing it
changes the drawing, not the finite-element resolution. Likewise,
|femm/arrow spacing| controls the displayed arrows. During a compilation,
further field pictures and profile plots reuse the model already solved;
a new compilation builds and solves it again. There is no persistent solver
cache between document runs.

\subsection{Measuring across both gaps}

Field lines show where the flux goes. Maya still needs numbers. She draws a
straight measurement contour halfway up the gaps:
\begin{codeexample}[code only]
\draw[femm/profile={name=coil-gap,samples=101},
  violet,dashed,thick,-Stealth] (-48,31.5)--(48,31.5);
\end{codeexample}
This path belongs inside the model picture. The |femm/profile| key records
it without creating another medium. It may come after the field picture,
since it does not change the material geometry. The chosen height, 31.5 mm,
lies inside the air gaps rather than exactly on an iron/air interface.
A measurement contour need not be closed and need not follow a field line.

After closing the model picture, Maya opens a new picture containing an axis:
\begin{codeexample}[code only]
\begin{tikzpicture}
\begin{axis}[width=13cm,height=5cm,grid=major,no markers,
  xlabel={$x$ (mm)},ylabel={Field (T)},
  legend style={at={(.5,1.03)},anchor=south,legend columns=3}]
  \femmplot[profile=coil-gap,component=B,abscissa=x,
    plot options={black,thick}]
  \addlegendentry{$|\mathbf B|$}
  \femmplot[profile=coil-gap,component=Bx,abscissa=x,
    plot options={blue}]
  \addlegendentry{$B_x$}
  \femmplot[profile=coil-gap,component=By,abscissa=x,
    plot options={red}]
  \addlegendentry{$B_y$}
\end{axis}
\end{tikzpicture}
\end{codeexample}
There is no |femm/problem| option on this graph. Each |\femmplot| names the
profile to sample; it already knows its source model. The command emits a
complete PGFPlots plot, including its terminating semicolon, so Maya adds
only the legend entry afterwards.

|component=B| means the magnitude; |Bx| and |By| retain their signs.
|abscissa=x| plots physical $x$ in millimetres. With the stated winding signs,
the dominant vertical component points upwards in the left gap and downwards
in the right one. The two gaps need not have exactly equal magnitudes:
the winding is on the left leg and the finite-domain solution includes leakage.

\clearpage
\subsection{The electromagnet, assembled}
\label{sec:tutorial-coil-result}
The complete file \nolinkurl{examples/tutorial-coil.tex} produces the following
picture and graph. These are computed with the draft settings just introduced. The small jumps in the curves come from the
piecewise-constant element field; PGFPlots joins the samples with line segments.
\begin{center}
  \includegraphics[width=.88\linewidth]{build/examples/tutorial-coil.pdf}
\end{center}

The triangle count and gap-centre values beneath the graph are printed by
|\femmstat{triangles}| and |\femmB|. They are useful for checking that a mesh
change actually took effect. The latter prints only three decimal places;
use profiles or the Lua API when more numerical precision is needed.

\clearpage
\subsection{Replacing the coil by a permanent magnet}

Maya would now like a field without supplying a current. She keeps the U,
removes \emph{both} copper rectangles, and inserts an alnico crosspiece at its
base. This time the armature will leave 1 mm gaps. She starts a new model
picture with the same outer air domain and the same 5 mm global mesh.

The |core| style remains unchanged. A new style supplies the hard material
and its magnetisation direction:
\begin{codeexample}[code only]
\tikzset{
  magnet/.style={femm/region={material=cast-alnico-5-lng37,
    magnetization angle=180},draw=red!70!black,fill=red!25},
  gap mesh/.style={femm/region={material=air,mesh size=.5}}
}
\end{codeexample}
After drawing the \emph{same U polygon}, she adds these paths in this order:
\begin{codeexample}[code only]
\filldraw[magnet] (-25,-30) rectangle (25,-15);
\filldraw[core] (-40,31) rectangle (40,46);
\path[gap mesh] (-40,30) rectangle (-25,31);
\path[gap mesh] (25,30) rectangle (40,31);
\end{codeexample}
The alnico rectangle overlaps the bottom of the original U. That is deliberate:
because it comes later, it replaces the iron in that part of the crosspiece.
The armature now extends from $y=31$ to $46$ mm. The gap rectangles must change
with it; leaving the old 3 mm refinement rectangles would paint air over part
of the new armature.

There are no current-density declarations in this model. The selected materials
have zero source current. The alnico record supplies the coercivity and shifted
B--H table, so there is no need to guess an equivalent coil current or enter a
second remanence value. The angle $180^\circ$ points the prescribed magnetisation
to the left. With this orientation, flux rises through the left leg, crosses the
armature, and returns through the right leg.

The material law describes a prescribed static branch. Changing the drawing
does not simulate the magnetic history of a real alnico magnet or irreversible
demagnetisation. Those limits still apply even when the computed field lines
look convincing.

\subsection{Keeping one contour and bending another}

Maya uses the same field-picture command as before. Her straight profile needs
a new name and a new height, halfway through the narrower gaps:
\begin{codeexample}[code only]
\draw[femm/profile={name=magnet-gap,samples=101},
  violet,dashed,thick,-Stealth] (-48,30.5)--(48,30.5);
\end{codeexample}
She copies the previous axis and replaces |profile=coil-gap| by
|profile=magnet-gap| in its three plots. Distinct names let both models and
both graphs coexist in one document. The two devices have different gaps and
different excitation; their field magnitudes are not expected to match.

She also wonders how the field varies along a curved contour in the central
air space. The profile key works on a Bezier path too:
\begin{codeexample}[code only]
\draw[femm/profile={name=magnet-arc,samples=101,
  curve tolerance=.03},orange!85!black,dashed,thick,-Stealth]
  (-18,-8) .. controls (-22,22) and (22,22) .. (18,-8);
\end{codeexample}
Again this is inside the model picture, after all physical regions. The
contour is a measurement path chosen by Maya, not a computed field line.
For its graph she chooses arc length and components relative to the path:
\begin{codeexample}[code only]
\begin{tikzpicture}
\begin{axis}[width=13cm,height=5cm,grid=major,no markers,
  xlabel={Arc length $s$ (mm)},ylabel={Field (T)},
  legend style={at={(.5,1.03)},anchor=south,legend columns=3}]
  \femmplot[profile=magnet-arc,component=B,
    plot options={black,thick}]
  \addlegendentry{$|\mathbf B|$}
  \femmplot[profile=magnet-arc,component=Bt,
    plot options={orange!85!black}]
  \addlegendentry{$B_t$}
  \femmplot[profile=magnet-arc,component=Bn,
    plot options={teal!70!black}]
  \addlegendentry{$B_n$}
\end{axis}
\end{tikzpicture}
\end{codeexample}
The default abscissa is |s|, the distance travelled along the flattened path,
in millimetres. It is not the Bezier parameter. |Bt| is the tangential component
and |Bn| is the component along the left normal. Reversing the traversal
reverses those component signs at a common point. Cartesian components remain
available by using |Bx| or |By| instead.

If Maya needs a closed measurement contour, she can use a path ending in
|-- cycle| with the same profile key. The first position then appears again
at the last sample. At a sharp corner the tangent changes, so local components
can jump even in a smooth physical field. On a material interface the sampling
side is ambiguous; placing the contour inside the desired medium or using a
small signed |offset| makes the choice explicit.

\clearpage
\subsection{The permanent magnet, assembled}
The file \nolinkurl{examples/tutorial-alnico.tex} combines the device and both
profiles. The orange curve lies in the central air space, whereas the violet
line crosses the pole gaps; they sample different parts of the same field.
\begin{center}
  \includegraphics[width=.82\linewidth]{build/examples/tutorial-alnico.pdf}
\end{center}

\clearpage
\subsection{Finishing the picture: now refine the computation}
\label{sec:tutorial-refinement}

Maya is happy with the geometry, labels and axes. Now it is worth spending more
time on the finite-element solution. She first refines the global spacing,
then the gaps, comparing results after each change instead of assuming that
one particular setting is sufficient.

\begin{center}
\begin{tabular}{llll}
\toprule Setting & Working draft & First refinement & Further gap refinement\\
\midrule
Global |mesh size| & 5 & 3 & 2\\
Coil gap |mesh size| & 1 & .5 & .25\\
Alnico gap |mesh size| & .5 & .25 & .125\\
\bottomrule
\end{tabular}
\end{center}
All entries are in millimetres. These are suggested trials, not accuracy
prescriptions or guaranteed mesh-quality bounds. For example, the first
refinement of the coil model consists of changing these two settings:
\begin{codeexample}[code only]
% In the model picture options:
femm/problem={xmin=-110,xmax=110,
  ymin=-95,ymax=115,mesh size=3}
% In the style definition:
gap mesh/.style={femm/region={material=air,mesh size=.5}}
\end{codeexample}
The physical geometry, material laws and current density stay fixed during
this comparison. Each trial must construct a new model, normally by recompiling
the document. Changing a key after |\pic {femm field};| does not remesh the
model that has already been solved.

She compares the field at fixed positions inside both gaps and the curves as
a whole. She looks especially at the places where the graph changes rapidly.
Then, keeping the mesh spacings fixed, she enlarges the outer air domain
by moving all four boundaries outwards. This separate check tells her whether
the artificial $A_z=0$ boundary still influences the displayed region.
A small Newton residual alone answers neither question.

For a direct overlay in one document, Maya can draw two complete model pictures,
using names such as |coil-coarse| and |coil-fine| for the same physical contour.
After both pictures, she plots them on one axis:
\begin{codeexample}[code only]
\femmplot[profile=coil-coarse,component=By,abscissa=x,
  plot options={gray,dashed}]
\addlegendentry{Coarse mesh}
\femmplot[profile=coil-fine,component=By,abscissa=x,
  plot options={red}]
\addlegendentry{Refined mesh}
\end{codeexample}
This fragment assumes both profiles have been registered under those names.
Each keeps its own solved model. Merely changing |samples=101| to |samples=301|
on the coarse model would sample the same approximate solution more densely;
it would not produce the fine-mesh comparison.

Once the values of interest are stable enough for her purpose, Maya can increase
the contour levels or profile samples for the final presentation. She keeps
the mesh/domain sensitivity in mind and retains the source settings alongside
the figure. The workflow has been economical: ordinary TikZ while arranging
the device, a coarse solve while developing the explanation, and more triangles
when there is a numerical result worth checking carefully.

The remaining sections provide the complete key reference and numerical
conventions. This tutorial follows the progressive, task-based presentation
of the PGF tutorials; its narrative and magnetic examples are original.
\clearpage
